% GENERATED FILE. Edit canonical lessons, then run npm run generate:books.
% canonical-source-hash: fnv1a64:b26491f52392d7b7
% canonical-lessons: KA-C333-neraverisu, KA-C333-samarambha, KA-C333-arpisu, KA-C333-erpadisu, KA-C333-amantrana, KA-R333-review-327-333

\chapter{Rites and Ceremonies}
\label{ch:ka-rites-and-ceremonies}
\hlchaptermodality{333}

\begin{chapteropening}
\textbf{By the end of this chapter:} \emph{I can say to carry out, a ceremony, to offer, to set up and an invitation card.}

Complete the last lesson of chapter 333: I can say to carry out, a ceremony, to offer, to set up and an invitation card.
\end{chapteropening}

\section[\texorpdfstring{\kn{ನೆರವೇರಿಸು} (neravērisu)}{ನೆರವೇರಿಸು (neravērisu)}]{\kn{ನೆರವೇರಿಸು} (neravērisu) — to carry out (a rite); to perform}

\label{lesson:KA-C333-neraverisu}



\begin{quote}
Before the new one: say the Kannada for a victim, then the Kannada for a forewarning.
\end{quote}

\subsection*{You'll want to know: \kn{ನೆರವೇರಿಸು}}
\textbf{\kn{ನೆರವೇರಿಸು}} — \emph{neravērisu} — \textquotedblleft{}to carry out (a rite); to perform\textquotedblright{}.

\begin{quote}
\kn{ಅರ್ಚಕರು} \kn{ಮದುವೆಯ} \kn{ಪೂಜೆಯನ್ನು} \kn{ನೆರವೇರಿಸಿದರು}.
\end{quote}

\begin{quote}
\emph{arcakaru maduveya pūjeyannu neravērisidaru.}
\end{quote}

\begin{quote}
The priest performed the wedding puja.
\end{quote}

Past (I …): \kn{ನೆರವೇರಿಸಿದೆ}, \emph{neravēriside}.

Rites and ceremonies are carried out with it. A wish or a plan that comes true \kn{ನೆರವೇರುತ್ತದೆ}. Before a rite is carried out, the house is made \kn{ಸ್ವಚ್ಛ}, spotless.

\begin{cousinweb}
\textbf{\kn{ನೆರವೇರಿಸು}} (\emph{neravērisu}) is usually traced to the native Dravidian \emph{neṟe}, \textquotedblleft{}to become complete, to be realized, to occur\textquotedblright{} (DEDR 3682), with the causative ending -\emph{isu}.
\end{cousinweb}

\begin{grammarlens}[title={What you've built}]
One more word for telling what happened.
\end{grammarlens}

\subsection*{Your turn}
\begin{itemize}
\raggedright
  \item \emph{Say it:} \emph{neravērisu}
  \item \emph{Say it:} \emph{neravērisu}, once more
  \item \emph{Say it:} \emph{arcakaru maduveya pūjeyannu neravērisidaru.} — The priest performed the wedding puja.
  \item \emph{Recall:} say the Kannada for a victim, then the Kannada for a forewarning, then say \emph{neravērisu} again
\end{itemize}

\begin{tcolorbox}[breakable,colback=teal!4,colframe=teal!35!black,title={Before you move on}]
What does \kn{ನೆರವೇರಿಸು} mean? (\textquotedblleft{}To carry out (a rite); to perform\textquotedblright{}.) Say it once more.
\end{tcolorbox}
\section[\texorpdfstring{\kn{ಸಮಾರಂಭ} (samārambha)}{ಸಮಾರಂಭ (samārambha)}]{\kn{ಸಮಾರಂಭ} (samārambha) — a ceremony; a function}

\label{lesson:KA-C333-samarambha}



\begin{quote}
Before the new one: say the Kannada for a forewarning, then the Kannada for to carry out.
\end{quote}

\subsection*{You'll want to know: \kn{ಸಮಾರಂಭ}}
\textbf{\kn{ಸಮಾರಂಭ}} — \emph{samārambha} — \textquotedblleft{}a ceremony; a function\textquotedblright{}.

\begin{quote}
\kn{ಶಾಲೆಯ} \kn{ವಾರ್ಷಿಕ} \kn{ಸಮಾರಂಭದಲ್ಲಿ} \kn{ಮಕ್ಕಳು} \kn{ಹಾಡಿದರು}.
\end{quote}

\begin{quote}
\emph{śāleya vārṣika samārambhadalli makkaḷu hāḍidaru.}
\end{quote}

\begin{quote}
The children sang at the school's annual function.
\end{quote}

Plural: \kn{ಸಮಾರಂಭಗಳು}, \emph{samārambhagaḷu}.

Any formal gathering: a wedding, a prize-giving, an opening.

\begin{cousinweb}
\textbf{\kn{ಸಮಾರಂಭ}} (\emph{samārambha}) is borrowed from Sanskrit \emph{samārambha}, \textquotedblleft{}an undertaking, an enterprise\textquotedblright{} (Monier-Williams); in Kannada it became an organised event.
\end{cousinweb}

\begin{grammarlens}[title={What you've built}]
One more word for telling what happened.
\end{grammarlens}

\subsection*{Your turn}
\begin{itemize}
\raggedright
  \item \emph{Say it:} \emph{samārambha}
  \item \emph{Say it:} \emph{samārambha}, once more
  \item \emph{Say it:} \emph{śāleya vārṣika samārambhadalli makkaḷu hāḍidaru.} — The children sang at the school's annual function.
  \item \emph{Recall:} say the Kannada for a forewarning, then the Kannada for to carry out, then say \emph{samārambha} again
\end{itemize}

\begin{tcolorbox}[breakable,colback=teal!4,colframe=teal!35!black,title={Before you move on}]
What does \kn{ಸಮಾರಂಭ} mean? (\textquotedblleft{}A ceremony; a function\textquotedblright{}.) Say it once more.
\end{tcolorbox}
\section[\texorpdfstring{\kn{ಅರ್ಪಿಸು} (arpisu)}{ಅರ್ಪಿಸು (arpisu)}]{\kn{ಅರ್ಪಿಸು} (arpisu) — to offer; to dedicate}

\label{lesson:KA-C333-arpisu}



\begin{quote}
Before the new one: say the Kannada for to carry out, then the Kannada for a ceremony.
\end{quote}

\subsection*{You'll want to know: \kn{ಅರ್ಪಿಸು}}
\textbf{\kn{ಅರ್ಪಿಸು}} — \emph{arpisu} — \textquotedblleft{}to offer; to dedicate\textquotedblright{}.

\begin{quote}
\kn{ಮಕ್ಕಳು} \kn{ಹಿರಿಯರಿಗೆ} \kn{ಹೂವು} \kn{ಅರ್ಪಿಸಿದರು}.
\end{quote}

\begin{quote}
\emph{makkaḷu hiriyarige hūvu arpisidaru.}
\end{quote}

\begin{quote}
The children offered flowers to the elders.
\end{quote}

Past (I …): \kn{ಅರ್ಪಿಸಿದೆ}, \emph{arpiside}.

What is offered takes -\kn{ಅನ್ನು} or stands bare; to whom takes -\kn{ಗೆ}. Books are dedicated too: \kn{ಪುಸ್ತಕವನ್ನು} \kn{ಅಮ್ಮನಿಗೆ} \kn{ಅರ್ಪಿಸು}.

\begin{cousinweb}
\textbf{\kn{ಅರ್ಪಿಸು}} (\emph{arpisu}) is built on Sanskrit \emph{arpaṇa}, \textquotedblleft{}procuring, placing upon, handing over\textquotedblright{} (Monier-Williams), with the Kannada verb ending -\emph{isu}.
\end{cousinweb}

\begin{grammarlens}[title={What you've built}]
One more word for telling what happened.
\end{grammarlens}

\subsection*{Your turn}
\begin{itemize}
\raggedright
  \item \emph{Say it:} \emph{arpisu}
  \item \emph{Say it:} \emph{arpisu}, once more
  \item \emph{Say it:} \emph{makkaḷu hiriyarige hūvu arpisidaru.} — The children offered flowers to the elders.
  \item \emph{Recall:} say the Kannada for to carry out, then the Kannada for a ceremony, then say \emph{arpisu} again
\end{itemize}

\begin{tcolorbox}[breakable,colback=teal!4,colframe=teal!35!black,title={Before you move on}]
What does \kn{ಅರ್ಪಿಸು} mean? (\textquotedblleft{}To offer; to dedicate\textquotedblright{}.) Say it once more.
\end{tcolorbox}
\section[\texorpdfstring{\kn{ಏರ್ಪಡಿಸು} (ērpaḍisu)}{ಏರ್ಪಡಿಸು (ērpaḍisu)}]{\kn{ಏರ್ಪಡಿಸು} (ērpaḍisu) — to set up; to arrange}

\label{lesson:KA-C333-erpadisu}



\begin{quote}
Before the new one: say the Kannada for a ceremony, then the Kannada for to offer.
\end{quote}

\subsection*{You'll want to know: \kn{ಏರ್ಪಡಿಸು}}
\textbf{\kn{ಏರ್ಪಡಿಸು}} — \emph{ērpaḍisu} — \textquotedblleft{}to set up; to arrange\textquotedblright{}.

\begin{quote}
\kn{ಅಪ್ಪ} \kn{ಮದುವೆಗೆ} \kn{ಊಟ} \kn{ಏರ್ಪಡಿಸಿದರು}.
\end{quote}

\begin{quote}
\emph{appa maduvege ūṭa ērpaḍisidaru.}
\end{quote}

\begin{quote}
Father arranged the meal for the wedding.
\end{quote}

Past (I …): \kn{ಏರ್ಪಡಿಸಿದೆ}, \emph{ērpaḍiside}.

\kn{ಏರ್ಪಡು} is to come about; \kn{ಏರ್ಪಡಿಸು} is to make it come about. The noun is \kn{ಏರ್ಪಾಡು}, an arrangement.

\begin{cousinweb}
\textbf{\kn{ಏರ್ಪಡಿಸು}} (\emph{ērpaḍisu}) is native Dravidian: DEDR 905 lists Kannada \emph{ērpaḍu}, \textquotedblleft{}to come about\textquotedblright{}, and \emph{ērpaḍisu}, beside Tamil \emph{ēl}, \textquotedblleft{}to happen, to be suitable\textquotedblright{}.
\end{cousinweb}

\begin{grammarlens}[title={What you've built}]
One more word for telling what happened.
\end{grammarlens}

\subsection*{Your turn}
\begin{itemize}
\raggedright
  \item \emph{Say it:} \emph{ērpaḍisu}
  \item \emph{Say it:} \emph{ērpaḍisu}, once more
  \item \emph{Say it:} \emph{appa maduvege ūṭa ērpaḍisidaru.} — Father arranged the meal for the wedding.
  \item \emph{Recall:} say the Kannada for a ceremony, then the Kannada for to offer, then say \emph{ērpaḍisu} again
\end{itemize}

\begin{tcolorbox}[breakable,colback=teal!4,colframe=teal!35!black,title={Before you move on}]
What does \kn{ಏರ್ಪಡಿಸು} mean? (\textquotedblleft{}To set up; to arrange\textquotedblright{}.) Say it once more.
\end{tcolorbox}
\section[\texorpdfstring{\kn{ಆಮಂತ್ರಣ} (āmantraṇa)}{ಆಮಂತ್ರಣ (āmantraṇa)}]{\kn{ಆಮಂತ್ರಣ} (āmantraṇa) — an invitation card; an invitation}

\label{lesson:KA-C333-amantrana}



\begin{quote}
Before the new one: say the Kannada for to offer, then the Kannada for to set up.
\end{quote}

\subsection*{You'll want to know: \kn{ಆಮಂತ್ರಣ}}
\textbf{\kn{ಆಮಂತ್ರಣ}} — \emph{āmantraṇa} — \textquotedblleft{}an invitation card; an invitation\textquotedblright{}.

\begin{quote}
\kn{ಮದುವೆಯ} \kn{ಆಮಂತ್ರಣ} \kn{ಅಂಚೆಯಲ್ಲಿ} \kn{ಬಂತು}.
\end{quote}

\begin{quote}
\emph{maduveya āmantraṇa añceyalli bantu.}
\end{quote}

\begin{quote}
The wedding invitation came by post.
\end{quote}

Plural: \kn{ಆಮಂತ್ರಣಗಳು}, \emph{āmantraṇagaḷu}.

In full the printed card is \kn{ಆಮಂತ್ರಣ} \kn{ಪತ್ರಿಕೆ}; \kn{ಆಹ್ವಾನ} is the general word for an invitation.

\begin{cousinweb}
\textbf{\kn{ಆಮಂತ್ರಣ}} (\emph{āmantraṇa}) is borrowed from Sanskrit \emph{āmantraṇa}, \textquotedblleft{}addressing, calling to, summoning\textquotedblright{} (Monier-Williams).
\end{cousinweb}

\begin{grammarlens}[title={What you've built}]
One more word for telling what happened.
\end{grammarlens}

\subsection*{Your turn}
\begin{itemize}
\raggedright
  \item \emph{Say it:} \emph{āmantraṇa}
  \item \emph{Say it:} \emph{āmantraṇa}, once more
  \item \emph{Say it:} \emph{maduveya āmantraṇa añceyalli bantu.} — The wedding invitation came by post.
  \item \emph{Recall:} say the Kannada for to offer, then the Kannada for to set up, then say \emph{āmantraṇa} again
\end{itemize}

\begin{tcolorbox}[breakable,colback=teal!4,colframe=teal!35!black,title={Before you move on}]
What does \kn{ಆಮಂತ್ರಣ} mean? (\textquotedblleft{}An invitation card; an invitation\textquotedblright{}.) Say it once more.
\end{tcolorbox}
\section[(dialogue)]{Looking back: from All of a Sudden to Rites and Ceremonies}

\label{lesson:KA-R333-review-327-333}



\begin{quote}
Say the Kannada for meanwhile and all of a sudden.
\end{quote}

\subsection*{Your turn}
Say these aloud:

\begin{itemize}
\raggedright
  \item meanwhile, all of a sudden, at once, abruptly, to be startled
  \item a robbery, to loot, to commit, to abscond, a track
  \item the media, a TV channel, a news bulletin, to announce, a statement
  \item to collide, an eyewitness, damage, the deceased, a mishap
  \item once before, on that day, a long time, ancient, an era
  \item a disaster, to be caught, to survive, a victim, a forewarning
  \item to carry out, a ceremony, to offer, to set up, an invitation card
\end{itemize}

\begin{tcolorbox}[breakable,colback=teal!4,colframe=teal!35!black,title={Before you move on}]
Meanwhile? (\textbf{\kn{ಅಷ್ಟರಲ್ಲಿ}}, \emph{aṣṭaralli}.) All of a sudden? (\textbf{\kn{ಹಠಾತ್ತಾಗಿ}}, \emph{haṭhāttāgi}.) At once? (\textbf{\kn{ಕೂಡಲೇ}}, \emph{kūḍalē}.) Abruptly? (\textbf{\kn{ಏಕಾಏಕಿ}}, \emph{ēkāēki}.) To be startled? (\textbf{\kn{ಬೆಚ್ಚಿಬೀಳು}}, \emph{beccibīḷu}.) A robbery? (\textbf{\kn{ದರೋಡೆ}}, \emph{darōḍe}.) To loot? (\textbf{\kn{ದೋಚು}}, \emph{dōcu}.) To commit? (\textbf{\kn{ಎಸಗು}}, \emph{esagu}.) To abscond? (\textbf{\kn{ಪರಾರಿಯಾಗು}}, \emph{parāriyāgu}.) A track? (\textbf{\kn{ಜಾಡು}}, \emph{jāḍu}.) The media? (\textbf{\kn{ಮಾಧ್ಯಮ}}, \emph{mādhyama}.) A TV channel? (\textbf{\kn{ವಾಹಿನಿ}}, \emph{vāhini}.) A news bulletin? (\textbf{\kn{ವಾರ್ತೆ}}, \emph{vārte}.) To announce? (\textbf{\kn{ಘೋಷಿಸು}}, \emph{ghōṣisu}.) A statement? (\textbf{\kn{ಹೇಳಿಕೆ}}, \emph{hēḷike}.) To collide? (\textbf{\kn{ಡಿಕ್ಕಿ} \kn{ಹೊಡೆ}}, \emph{ḍikki hoḍe}.) An eyewitness? (\textbf{\kn{ಪ್ರತ್ಯಕ್ಷದರ್ಶಿ}}, \emph{pratyakṣadarśi}.) Damage? (\textbf{\kn{ಹಾನಿ}}, \emph{hāni}.) The deceased? (\textbf{\kn{ಮೃತ}}, \emph{mṛta}.) A mishap? (\textbf{\kn{ಅವಘಡ}}, \emph{avaghaḍa}.) Once before? (\textbf{\kn{ಹಿಂದೊಮ್ಮೆ}}, \emph{hindomme}.) On that day? (\textbf{\kn{ಅಂದು}}, \emph{andu}.) A long time? (\textbf{\kn{ಬಹುಕಾಲ}}, \emph{bahukāla}.) Ancient? (\textbf{\kn{ಪ್ರಾಚೀನ}}, \emph{prācīna}.) An era? (\textbf{\kn{ಯುಗ}}, \emph{yuga}.) A disaster? (\textbf{\kn{ವಿಪತ್ತು}}, \emph{vipattu}.) To be caught? (\textbf{\kn{ಸಿಲುಕು}}, \emph{siluku}.) To survive? (\textbf{\kn{ಬದುಕುಳಿ}}, \emph{badukuḷi}.) A victim? (\textbf{\kn{ಸಂತ್ರಸ್ತ}}, \emph{santrasta}.) A forewarning? (\textbf{\kn{ಮುನ್ಸೂಚನೆ}}, \emph{munsūcane}.) To carry out? (\textbf{\kn{ನೆರವೇರಿಸು}}, \emph{neravērisu}.) A ceremony? (\textbf{\kn{ಸಮಾರಂಭ}}, \emph{samārambha}.) To offer? (\textbf{\kn{ಅರ್ಪಿಸು}}, \emph{arpisu}.) To set up? (\textbf{\kn{ಏರ್ಪಡಿಸು}}, \emph{ērpaḍisu}.) An invitation card? (\textbf{\kn{ಆಮಂತ್ರಣ}}, \emph{āmantraṇa}.)
\end{tcolorbox}
